\documentclass[11pt]{article}
\usepackage[margin=1in]{geometry}
\usepackage{amsmath,amssymb,amsthm,bm,booktabs,graphicx,hyperref,enumitem,xcolor}
\usepackage[T1]{fontenc}
\hypersetup{colorlinks=true,linkcolor=blue!50!black,urlcolor=blue!50!black}
\setlist{itemsep=2pt,topsep=3pt}
\newcommand{\R}{\mathbb{R}}
\newcommand{\E}{\mathbb{E}}
\newcommand{\Var}{\operatorname{Var}}
\newcommand{\Cov}{\operatorname{Cov}}
\newcommand{\corr}{\operatorname{corr}}
\newcommand{\tr}{\operatorname{tr}}
\newcommand{\orth}{\operatorname{orth}}
\newcommand{\med}{\operatorname{med}}
\newcommand{\rank}{\operatorname{rank}}
\newcommand{\Rc}{R^2_{\mathrm{c}}}
\newtheorem*{remark}{Remark}
\title{Predictive neural subspaces: the whole project in mathematical terms}
\author{Written for Gunnar Andra-Thomas \\ \small Branch \texttt{pilot-timescale-fix}, state of 24 November 2026}
\date{}
\begin{document}
\maketitle
\tableofcontents

\section*{How to read this}

Every analysis in the project is written down here as the exact quantity the code computes, with the file that
computes it, and then every result is stated and interpreted in those terms. Part I is the original, pre-registered
study; Part II the post hoc follow-ups after the talk (A17); Part III the calibrated simulation and the latent model
(A18). Section~\ref{sec:synthesis} pulls the mathematics of the failure together. Numbers are mouse means over 13
mice unless they are called session medians (over 16 sessions) or seed means (over simulated sessions).

%======================================================================
\part{The registered study}
%======================================================================

\section{Data and preprocessing}

\paragraph{Sessions, units, bins.} Session $\sigma$ (16 sessions, 13 mice; $m(\sigma)$ is its mouse) has $p$ units
(22 to 128, median 48.5) that fire at least $0.1$ Hz over the task epoch. The epoch is cut into $T$ bins of
$\delta = 0.5$ s (median $T \approx 7{,}980$, about 66 min). Let $n_{tu}$ be the spike count of unit $u$ in bin $t$.

\paragraph{Variance-stabilised, z-scored activity.} With $s_{tu} = \sqrt{n_{tu}}$,
\begin{equation}
z_{tu} = \frac{s_{tu} - \mu_u}{\max(\hat\sigma_u, \sigma_{\min})},
\qquad \sigma_{\min} = \operatorname{SD}\!\big[\sqrt{N}\big],\ N \sim \mathrm{Poisson}(0.1 \cdot 0.5) \approx 0.218,
\label{eq:z}
\end{equation}
where $\mu_u, \hat\sigma_u$ are fitted on the bins a model is fitted on and then reused unchanged for everything that
model touches (the D10 fix: the floor stops a unit silent in one fold from exploding in another;
\texttt{transform.py}). The population vector is $\bm z_t \in \R^p$.

\paragraph{Trial windows.} For trial $k$ with stimulus onset $t^{\mathrm{on}}_k$, the pre-stimulus window is
$[t^{\mathrm{on}}_k - 0.6,\ t^{\mathrm{on}}_k - 0.1)$ s; $\bm w_k \in \R^p$ is its square-rooted count vector,
z-scored with the session fit's $\mu,\hat\sigma$. Post-stimulus, choice and feedback windows are defined the same way.
The decoder uses the trials that pass the IBL Brain-Wide-Map mask and have a finite window (413 to 700 per session).

\section{The task and the Bayes-optimal prior}

\paragraph{Generative process.} Stimulus side $c_k \in \{+1, -1\}$ (left, right). The session starts with 90
unbiased trials ($P(c_k = +1) = 0.5$), then alternates blocks with $P(c_k = +1) = \beta \in \{0.2, 0.8\}$ whose lengths
are drawn from an exponential with mean 60, truncated to $(20, 100)$ trials.

\paragraph{Ideal observer.} The hidden state is $(\ell_k, \beta_k)$: the number of trials into the current block and
the block type. With hazard $h(\ell) \propto e^{-\ell/60}\,\mathbf 1[20 \le \ell \le 100]$ (normalised over the
remaining support), a forward filter over the $100 \times 3$ states gives the predictive distribution
$P(\beta_k \mid c_{1:k-1})$ and the \emph{block prior}
\begin{equation}
\pi_k \;=\; P(c_k = +1 \mid c_{1:k-1}) \;=\; \sum_{\beta} P(\beta_k = \beta \mid c_{1:k-1})\, \beta,
\label{eq:prior}
\end{equation}
a deterministic function $\pi_k = f(c_1, \dots, c_{k-1})$ of the \emph{past} stimuli only (\texttt{prior.py}, a port of
the IBL \texttt{optimal\_Bayesian}). The last stimulus alone explains $R^2 = 0.42$ of $\Var(\pi)$
($\corr(\pi_k, c_{k-1}) = 0.64$), the last five explain $0.81$ and the last ten $0.86$ (session medians,
recorded tasks). This fact returns in Part III.

\paragraph{Pseudo sessions.} The null redraws the whole block structure from the generative process:
$\tilde c^{(b)}_{1:n}$, $b = 1, \dots, 200$, independent of the real $c$, and $\tilde\pi^{(b)} = f(\tilde c^{(b)})$.

\section{Lagged reduced-rank ridge regression}

\paragraph{Pairs and objective.} For lag $\Delta$ (in bins; registered $\Delta \in \{1, 4, 16, 60\}$, i.e.\ 0.5, 2,
8, 30 s), stack valid pairs $(t, t+\Delta)$ into $X = [\bm z_t^\top]$, $Y = [\bm z_{t+\Delta}^\top] \in \R^{N\times p}$
(centred). Reduced-rank ridge regression solves
\begin{equation}
\hat B_\Delta \;=\; \arg\min_{\rank B \le d}\ \|Y - XB\|_F^2 + \lambda \|B\|_F^2,
\qquad \lambda = \alpha\, \tr(X^\top X)/p, \quad d = 5.
\label{eq:rrr}
\end{equation}
Its solution (Mukherjee \& Zhu 2011; \texttt{rrr.py}) is built from the full-rank ridge solution
$B_R = (X^\top X + \lambda I)^{-1} X^\top Y$ and the $p\times p$ matrix
\begin{equation}
G \;=\; B_R^\top X^\top X B_R + \lambda B_R^\top B_R \;=\; \hat Y^\top \hat Y + \lambda B_R^\top B_R,
\qquad G = V \Lambda V^\top,\ \lambda_1 \ge \dots \ge \lambda_p,
\label{eq:G}
\end{equation}
as $\hat B_\Delta = B_R V_d V_d^\top$ with $V_d$ the top $d$ eigenvectors. Two bases follow:
\begin{align}
V^{\mathrm{out}}_\Delta &= V_d && \text{(loading basis: where the prediction lives)},\\
V^{\mathrm{in}}_\Delta &= \orth(B_R V_d) && \text{(readout basis: which current directions are used)}.
\end{align}
The key structural fact is in \eqref{eq:G}: $V^{\mathrm{out}}$ consists of the directions of \emph{largest predicted
variance} of $\bm z_{t+\Delta}$. Nothing in \eqref{eq:rrr} refers to the prior; the objective is unsupervised and
variance-ranked.

\paragraph{Choosing $\alpha$.} Over the grid $\alpha \in \{10^{-3}, \dots, 10^{3}\}$ (13 values), the held-out pooled
$R^2$ of predicting $\bm z_{t+\Delta}$ is computed with 5 contiguous folds over pairs, purging any training pair within
10 bins of a test block, z-scoring within training bins. One $\alpha$ per session maximises the mean of these curves
over the registered lags. Comparison subspaces: PCA (top $d$ principal directions of $Z$) and covariance-matched random
subspaces $\orth(E \Lambda_C^{1/2} \Gamma)$, $\Gamma_{ij} \sim \mathcal N(0,1)$, with $C = E\Lambda_C E^\top$ the
covariance of $Z$ (Elsayed et al.\ 2016).

\section{Decoding, nulls and tests}

\paragraph{Decoder.} For a basis $V \in \R^{p \times d}$ (or $V = I_p$, the full population), features are
$\bm x_k = V^\top \bm w_k$. A ridge regression with penalty 1 on standardised features is fitted on 4 of 5 contiguous
blocks of trials and predicts the fifth; the score is the pooled out-of-fold
\begin{equation}
R^2(y; V) \;=\; 1 - \frac{\sum_f \sum_{k \in \mathcal T_f} (y_k - \hat y_k)^2}{\sum_f \sum_{k\in\mathcal T_f} (y_k - \bar y_{\mathcal S_f})^2},
\end{equation}
with $\bar y_{\mathcal S_f}$ the training mean of fold $f$. Binary targets use L2 logistic regression and the mean
per-fold AUC.

\paragraph{Pseudo-session correction.} For the prior,
\begin{equation}
\Rc(V) \;=\; R^2(\pi; V) \;-\; \med_{b}\, R^2\big(\tilde\pi^{(b)}; V\big).
\label{eq:Rc}
\end{equation}
Contiguous folds make raw $R^2$ of a slow target strongly negative (extrapolating a block structure across folds);
the correction removes what any block-like sequence would score. The same 200 pseudo sequences are used for every
basis and lag, so differences get a paired null.

\paragraph{The three primary tests.} Let $\bar{(\cdot)}$ denote averaging sessions within mouse; tests are one-sided
Wilcoxon signed-rank over the 13 mouse means, Holm-corrected over the three.
\begin{itemize}
\item \textbf{P1 (geometry).} Principal angles $\theta_j = \arccos s_j$, with $s_j$ the singular values of
  $\orth(A)^\top \orth(B)$; summary $\bar\theta = \frac1d \sum_j \theta_j$. With readout bases fitted on the two
  session halves, $D_{\mathrm{short}} = \bar\theta(\text{cross}) - \bar\theta(0.5\text{ s}, 0.5\text{ s})$ and
  $D_{\mathrm{long}} = \bar\theta(\text{cross}) - \bar\theta(30\text{ s}, 30\text{ s})$; P1 is the intersection-union
  test $p = \max(p_{\mathrm{short}}, p_{\mathrm{long}})$.
\item \textbf{P2 (prior).} $\Delta_\pi = [R^2(\pi; V^{\mathrm{in}}_{60}) - R^2(\pi; V^{\mathrm{in}}_{1})] -
  \med_b[R^2(\tilde\pi^{(b)}; V^{\mathrm{in}}_{60}) - R^2(\tilde\pi^{(b)}; V^{\mathrm{in}}_{1})]$; test $\bar\Delta_\pi > 0$.
\item \textbf{P3 (current stimulus).} The same with corrected AUC of $c_k$ on the post-stimulus window, short minus long.
\end{itemize}

\paragraph{Elapsed time.} The bin time $t$ is decoded from $V^\top \bm z_t$ with 20 contiguous chunks assigned
round-robin to 5 folds; a high $R^2_{\mathrm{el}}$ means the subspace's activity varies smoothly with session time.

\section{Results of the registered study}

\begin{center}
\begin{tabular}{lll}
\toprule
Test & Statistic (mouse mean [95\% CI]) & Holm $p$ \\
\midrule
P1 & $\bar D_{\mathrm{short}} = 12.7^\circ\,[10.2, 15.1]$, $\bar D_{\mathrm{long}} = 3.6^\circ\,[2.2, 4.9]$ & 0.0026 \\
P2 & $\bar\Delta_\pi = -0.099\,[-0.218, 0.002]$ & 0.89 \\
P3 & $+0.061\,[0.028, 0.094]$ & 0.0061 \\
\bottomrule
\end{tabular}
\end{center}

Per basis, $\Rc$ is $0.054$ ($V^{\mathrm{in}}_1$), $-0.027$, $-0.057$, $-0.046$ ($V^{\mathrm{in}}_{60}$), $0.013$
(PCA), $-0.030$ (random); the full population scores $\Rc(I_p) = 0.159$ (session median 0.108). So the prior is
decodable from the population but not concentrated in any 5-dimensional summary, and P2 fails. What every summary
does contain is session time: $R^2_{\mathrm{el}} = 0.84$ ($V^{\mathrm{in}}_{60}$), $0.82$ ($V^{\mathrm{in}}_1$),
$0.80$ (PCA), $0.64$ (random).

\begin{remark}
P1 holds mostly through $D_{\mathrm{short}}$; the 30 s split-half angle is itself at the covariance-matched chance
level, so $D_{\mathrm{long}}$ is measured against noise. P3 holds because long-lag bases lose the stimulus (evoked
responses decay within a second), not because the 0.5 s basis is specialised for it: its corrected AUC equals the
random basis's.
\end{remark}

%======================================================================
\part{Why the objective could not be expected to find the prior (A17, post hoc)}
%======================================================================

\section{The prior's timescale}

\paragraph{Definitions.} Assign $\pi_k$ to the bins between stimulus onsets $(t^{\mathrm{on}}_{k-1}, t^{\mathrm{on}}_k]$
to get a per-bin series $\pi(t)$, and define
\begin{equation}
\rho_\pi(\Delta) = \corr\big(\pi(t), \pi(t+\Delta)\big), \qquad
\tau_\pi = \delta\Big(\sum_{j=0}^{j_0 - 1} \rho_\pi(j\delta) - \tfrac12\Big),
\end{equation}
with $j_0$ the first lag where $\rho_\pi \le 0$ (for $\rho = e^{-t/\tau}$ this returns $\approx \tau$). The
trial-indexed version uses $\pi_k$ directly with lags in trials. $\rho_\pi(\Delta)^2$ is the fraction of $\Var(\pi)$
that is linearly predictable $\Delta$ ahead from the current value.

\paragraph{Results (task data only; session medians).}
\begin{center}
\begin{tabular}{lccccccc}
\toprule
$\Delta$ & 0.5 s & 2 s & 8 s & 30 s & 60 s & 2 min & 4 min \\
$\rho_\pi(\Delta)^2$ & 0.99 & 0.97 & 0.88 & 0.59 & 0.27 & 0.01 & 0.04 \\
\bottomrule
\end{tabular}
\end{center}
$\tau_\pi = 67$ s (IQR 55 to 75 s); trial-indexed $\tau_\pi = 12.2$ trials (IQR 10.9 to 13.8), first zero crossing at
27 trials. Blocks last 46 trials (3.6 min); sessions 66.5 min. The per-bin timescale is longer in trials (16.9)
because $\pi$ is constant through long inter-trial gaps.

\section{Variance shares}

On whole-session $Z$ (rows: bins), with $P^\perp_D$ the projection off $[\bm 1, D]$:
\begin{equation}
S_{\mathrm{low}} = 1 - \frac{\|P^\perp_{D_{20}} Z\|_F^2}{\|Z - \bar Z\|_F^2},
\qquad
S_\pi = \frac{\sum_u (\bm x^\top \bm w_u)^2 / (\bm x^\top \bm x)}{\sum_u \bm w_u^\top \bm w_u},
\end{equation}
where $D_{20}$ is the cosine basis $\cos(\pi j u)$, $u \in [0,1]$ the normalised session time, $j \le \lfloor 2L/20
\rfloor$ ($L$ in minutes: periods of at least 20 min), and for $S_\pi$ the columns $\bm w_u$ (activity) and $\bm x$
(prior) are first residualised on the same nuisance design. $S_\pi$ is corrected as $S_\pi(\pi) - \med_b S_\pi(\tilde
\pi^{(b)})$. Session medians: $S_{\mathrm{low}} = 22.6\%$, trial-locked share (A11 GLM) $8.6\%$, $S_\pi = 0.20\%$
(pre-stimulus windows; positive in 14 of 16) or $0.09\%$ (per bin). These are measured on different data, so they are
scales, not one additive budget: ``about two orders of magnitude''.

\section{The variance-ranking heuristic, and why it is not a theorem}

Suppose activity were a sum of independent components, $\bm z_t = \sum_j \bm a_j x_j(t) + \bm\epsilon_t$, each
$x_j$ with variance $\sigma_j^2$ and autocorrelation $\rho_j$. If the predictor could read each component perfectly, the
predictable variance it contributes at lag $\Delta$ is
\begin{equation}
\Var\big(\E[\bm a_j x_j(t+\Delta) \mid x_j(t)]\big) = \|\bm a_j\|^2 \sigma_j^2 \rho_j(\Delta)^2 .
\label{eq:heur}
\end{equation}
For a ramp-like drift $\rho_D(\Delta) \approx 1$ at every lag tested, while $\rho_\pi(\Delta)^2$ falls from 0.99 to
0.01; and $\sigma_\pi^2 \ll \sigma_D^2$. Under \eqref{eq:heur} the prior's share of predictable variance relative to
the drift therefore only falls with lag: longer lags should not help. A17 compared the prior's term
$\widehat{SS}_\pi\, \rho_\pi(\Delta)^2$ with $\lambda_5(G)$ (median ratio 0.05 to 0.66; below 1 in 16, 16, 16 and 12
of 16 sessions for raw/detrended at 0.5 s/30 s).

\paragraph{Why this is not an upper bound (withdrawn claim).} It was first called an upper bound implying that the
prior ``could not'' enter $V_d$. Four reasons it is not:
\begin{enumerate}
\item The best predictor of $\pi(t+\Delta)$ uses all of the past, $\E[\pi(t+\Delta) \mid \mathcal F_t]$, whose explained
  fraction can exceed $\rho_\pi(\Delta)^2$ (e.g.\ the run length $\ell_k$ predicts a switch).
\item Components are not separable in $G$: for a direction $\bm v = \bm v_\pi + \bm v_D$,
  $\bm v^\top G \bm v$ contains cross terms, so a weak component can tilt $V_d$ without its own term exceeding
  $\lambda_5$.
\item $G$ ranks \emph{output} directions; the prior is decoded through the \emph{input} basis
  $V^{\mathrm{in}} = \orth(B_R V_d)$.
\item $\widehat{SS}_\pi$ is a noisy, null-subtracted estimate, not a decomposition.
\end{enumerate}
Equation \eqref{eq:heur} is a guide to intuition, and Part III tests the intuition directly.

\section{Lag sweep to 4 minutes on drift-removed data}

Activity was detrended per unit, $Z \mapsto P^\perp_{D} Z$ with $D = D_{20}$ (cosine) or $D = [u, u^2]$ (quadratic),
and RRR refitted for $\Delta \in \{1, 4, 16, 60, 120, 240, 480\}$ bins with $\alpha$ fixed to the A16 value. A
reproduction check matched the original lags to $2\times10^{-14}$. Committed reading: ``band selection'' iff, in
both copies, $\Rc(V^{\mathrm{in}}_{120}) > \Rc(V^{\mathrm{in}}_{1})$ and $> \Rc(\mathrm{PCA})$ (paired null,
one-sided Wilcoxon, $p < 0.05$).

\begin{center}
\begin{tabular}{lccccccc}
\toprule
$\Delta$ & 0.5 s & 2 s & 8 s & 30 s & 60 s & 2 min & 4 min \\
\midrule
$\rho_\pi^2$ & 0.99 & 0.97 & 0.88 & 0.59 & 0.27 & 0.01 & 0.04 \\
$\Rc$, cosine copy & 0.031 & 0.043 & 0.037 & 0.051 & 0.049 & 0.032 & 0.038 \\
$\Rc$, quadratic copy & $-0.002$ & $-0.054$ & $-0.003$ & $-0.015$ & $-0.018$ & $-0.010$ & 0.008 \\
\bottomrule
\end{tabular}
\end{center}
All four tests fail ($p = 0.29, 0.12$ cosine; $0.55, 0.53$ quadratic): ``no band selection''. The committed
prediction $\Rc(\Delta) \propto \rho_\pi(\Delta)^2$ fails too: Spearman correlation $-0.14$ ($p = 0.76$). The prior
content is flat in $\Delta$ even where $\rho_\pi^2 \approx 0$, which is what one expects if the small $\Rc$ is
incidental to any 5-dimensional summary (PCA 0.028, random 0.020) rather than selected by predictability.

%======================================================================
\part{Testing the mechanism where the truth is known (A18, post hoc)}
%======================================================================

\section{The generative model}

Each simulated session copies one recorded session's task exactly (all $t^{\mathrm{on}}_k$, movement and feedback
times, $c_k$, contrasts, choices, outcomes, block probabilities, the trial mask; hence the true $\pi_k$), its $p$ and
each unit's mean rate $\bar r_u$. On a 50 ms grid, with $k(t)$ the trial whose task latents are active at $t$
(switching at feedback$_{k-1} + 0.5$ s, never later than $t^{\mathrm{on}}_k - 0.6$ s),
\begin{align}
\log \lambda_u(t) \;=\;& b_u
 + g_\pi\, \ell^{\pi}_u \big(\pi_{k(t)} - \tfrac12\big)
 + g_h\, \ell^{h}_u\, h_{k(t)}
 + g_D \textstyle\sum_{r=1}^{3} \ell^{D}_{ur} D_r(t)
 + g_{\mathrm{ev}} \textstyle\sum_{e} \ell^{e}_u E_e(t) \nonumber\\
& + \ell^{f}_u \big(g_s\, \textstyle\sum_k c_k \sqrt{\kappa_k}\, \phi(t - t^{\mathrm{on}}_k) + a_f(t)\big)
 + \ell^{a}_u a_{s}(t) + \bm m_u^\top \bm\eta(t),
\label{eq:gen}\\
n_{tu} \;\sim\;& \mathrm{Poisson}\big(\lambda_u(t)\, \cdot 0.05\big), \nonumber
\end{align}
where:
\begin{itemize}
\item $\phi(s) = e^{-s/0.3}\mathbf 1[0 \le s < 1.5]$ (seconds); $E_e(t) = \sum_k \phi(t - t^e_k)$ for $e \in$
  \{stimulus, movement, feedback\}: unsigned evoked responses; $\kappa_k$ is the contrast;
\item $h_k = \mathrm{std}\big(\sum_{j=1}^{10} c_{k-j}\, e^{-(j-1)/\tau_h}\big)$ with $\tau_h = 1$ trial: a stimulus
  memory that carries no block information except through the stimuli themselves;
\item $D_r(t)$, the drift processes: a random-sign ramp plus cosines with periods $\ge 20$ min (amplitudes $\propto
  1/j$), plus weight $g_{\mathrm{mid}}$ on a 5 to 20 min mixture, unit variance, variances decaying by 0.3 per $r$
  (``low-rank''); or one such process per unit (``per-unit'');
\item $a_s$ an AR(1) slow noise (autocorrelation 0.99 per 0.5 s), $a_f$ a fast AR(1), $\bm\eta$ two shared fast noise
  factors;
\item loadings $\ell$ i.i.d.\ Gaussian (scale 0.6), with the stimulus loading sharing 0.3 of the prior loading;
\item $b_u$ set after all terms so that the session mean of $\lambda_u$ equals $\bar r_u$.
\end{itemize}
Then the \emph{registered} pipeline is applied unchanged.

\paragraph{Calibration.} Six gains $\theta = (g_\pi, g_D, g_{\mathrm{mid}}, g_h, g_{\mathrm{ev}}, g_s)$ are matched to
six recorded session medians $m^\star$ by the damped fixed point
\begin{equation}
\theta_i \leftarrow \theta_i \cdot \operatorname{clip}\!\Big[\big(m^\star_i / m_i(\theta)\big)^{0.7 e_i}, \tfrac12, 2\Big],
\end{equation}
$e_i = \frac12$ for shares and $R^2$ (which scale like $\theta^2$), $e_i = 1$ for $\mathrm{AUC} - \frac12$, with
$m_i(\theta)$ the median over 16 simulated sessions (seeds 900 to 915). Achieved / target: full-population
$\Rc$ $0.111/0.108$; $S_{\mathrm{low}}$ $22.6\%/22.6\%$; 5 to 20 min share $4.7\%/4.7\%$; trial-locked $8.6\%/8.6\%$;
previous-side AUC $0.558/0.563$; current-side AUC $0.866/0.869$. Calibrated $g_\pi = 0.29$ and $g_h = 0.036$: the
prior alone already makes $c_{k-1}$ as decodable as in the recordings, because $\corr(\pi_k, c_{k-1}) = 0.64$.

\section{Why stimulus memory leaks through the pseudo-session null}
\label{sec:leak}

Suppose the population encodes no prior, only $h_k$, a linear function of $c_{k-1}, \dots, c_{k-10}$. The decoder
learns $\hat\pi_k = \bm\beta^\top \bm w_k \approx \gamma h_k$. For the real target, $\Cov(h_k, \pi_k) \neq 0$ because both
are functions of the same past stimuli; the out-of-fold $R^2(\pi)$ approaches $\corr(h, \pi)^2 > 0$ (minus estimation
loss). For each pseudo target, $\tilde\pi^{(b)}$ is built from $\tilde c^{(b)}$, independent of $c$, so
$\Cov(h_k, \tilde\pi^{(b)}_k) = 0$ and $R^2(\tilde\pi^{(b)}) \le 0$ in expectation. Hence
\begin{equation}
\E[\Rc] \;\approx\; \corr(h, \pi)^2 - (\text{estimation loss}) \;>\; 0 \quad\text{with no prior encoded.}
\end{equation}
The pseudo-session null is a valid null for ``activity unrelated to the task sequence'', not for ``activity
related to the stimuli but not to the belief''.

\paragraph{The history-controlled score.} Replace every target by its residual on its own recent stimuli,
\begin{equation}
\pi^{\perp}_k = \pi_k - P_{[\bm 1, c_{k-1}, \dots, c_{k-5}]}\,\pi_k, \qquad
\tilde\pi^{(b)\perp}_k = \tilde\pi^{(b)}_k - P_{[\bm 1, \tilde c^{(b)}_{k-1}, \dots, \tilde c^{(b)}_{k-5}]}\,\tilde\pi^{(b)}_k,
\end{equation}
and compute $\Rc$ on these. If the population encodes only $h$ with memory inside the 5-trial window, $\pi^\perp$ is
(nearly) orthogonal to what it encodes and $\E[R^{2,\perp}_{\mathrm{c}}] \approx 0$. But the last five stimuli explain $R^2 = 0.81$
of $\Var(\pi)$, so $\Var(\pi^\perp) \approx 0.19\,\Var(\pi)$: the control also removes most of a genuine prior. A memory
longer than 5 trials would both mimic $\pi$ more and leak past the control.

\section{Simulation design and readings}

Eleven conditions, 32 seeds each (seeds $s$ and $s+16$ share a recorded task, so 16 independent tasks), sharing task,
latent noise, most loadings, pseudo draws and fit seed across conditions. Per condition, with $\bar{\cdot}$ the seed
mean, the \emph{capture} of a basis is $\kappa(V) = \overline{\Rc(V)} / \overline{\Rc(I_p)}$. Pre-set readings
(committed before any simulated result):
\begin{itemize}
\item G1: positive control finds the prior ($\Rc(I_p) > 0$ and $\Rc(V^{\mathrm{in}}_{60}) > 0$ in $\ge 22$ of 32).
\item A: with $\kappa_0 = \kappa(V^{\mathrm{in}}_{60})$ without drift: ``weak prior sufficient'' if $\kappa_0 < 0.25$;
  else ``drift sufficient'' if $\kappa(V^{\mathrm{in}}_{60}) < \kappa_0/2$ and $\med\,\Delta_\pi \le 0$ at calibrated
  drift; ``not sufficient'' if $\ge \kappa_0/2$ and $> 0$; else ``intermediate''.
\item B: null calibrated; stimulus memory at calibrated strength leaks. B2 (added after calibration): memory alone
  tuned to the recorded previous-side AUC mimics the prior ($\ge 22$ of 32 positive and mean $\ge 0.054$).
\item C: 8 profile features compared with the recordings' 95\% intervals (descriptive).
\end{itemize}

\section{Simulation results}

\begin{center}
\begin{tabular}{lcccc}
\toprule
Condition & $\Rc(I_p)$ & $\Rc(V^{\mathrm{in}}_{60})$ & $\Rc(V^{\mathrm{in}}_{1})$ & $\Delta_\pi$ \\
\midrule
strong prior ($3g_\pi$), no drift & 0.632 & 0.537 & 0.449 & $+0.089$ \\
calibrated prior, no drift & 0.170 & 0.069 & 0.019 & $+0.050$ \\
low-rank drift $\times\frac13$ & 0.174 & 0.047 & 0.019 & $+0.027$ \\
low-rank drift $\times 1$ & 0.167 & $-0.006$ & $-0.033$ & $+0.027$ \\
low-rank drift $\times 3$ & 0.268 & $-0.205$ & $-0.172$ & $-0.036$ \\
per-unit drift $\times 1$ & 0.063 & $-0.088$ & $-0.038$ & $-0.052$ \\
per-unit drift $\times 3$ & 0.084 & 0.136 & 0.028 & $+0.097$ \\
no prior, no memory & $-0.043$ & $-0.036$ & $-0.053$ & $+0.017$ \\
no prior, memory alone matching recordings & 0.078 & $-0.010$ & $-0.023$ & $+0.016$ \\
\midrule
recordings & 0.159 & $-0.046$ & 0.054 & $-0.099$ \\
\bottomrule
\end{tabular}
\end{center}

\paragraph{Reading the table in terms of \eqref{eq:rrr}--\eqref{eq:G}.}
\begin{enumerate}
\item \emph{G1 passes} ($32/32$): when $g_\pi$ is large the prior direction carries enough predictable variance to
  enter $V_d$, and $\kappa(V^{\mathrm{in}}_{60}) = 0.85$.
\item \emph{Without drift, long lags favour the prior}: $\kappa(V^{\mathrm{in}}_{60}) = 0.41\,[0.33, 0.48]$ and
  $\Delta_\pi = +0.050$. With only fast competitors (stimulus and evoked responses, $\rho \approx 0$ beyond a second),
  \eqref{eq:heur} favours the slow prior at 30 s, as the original hypothesis assumed.
\item \emph{Calibrated drift erases the 30 s prior}: $\Rc(V^{\mathrm{in}}_{60})$ falls from 0.069 to $-0.006$ (paired
  drop positive in 23 of 32; post hoc Wilcoxon $p = 0.002$) while $\Rc(I_p)$ is unchanged (0.167): the information is
  still in the population, but $V_d$ is spent on the drift's large, persistent predicted variance. But $\Delta_\pi$
  stays positive (21 of 32 seeds), unlike the recordings' $-0.099$, so reading A is \textbf{intermediate}: drift
  reproduces the loss at 30 s, not the recordings' advantage for 0.5 s.
\item \emph{Profile}: at calibrated low-rank drift, 6 of 8 features fall inside the recordings' intervals, notably
  $R^2_{\mathrm{el}} = 0.79, 0.75, 0.65$ for $V^{\mathrm{in}}_{60}$, PCA, random (recordings $0.84, 0.80, 0.64$).
  Per-unit drift matches 2 of 8 ($R^2_{\mathrm{el}} \approx 0.92$, too high) and makes even $\Rc(I_p)$ unreliable.
\item \emph{Stimulus memory mimics the prior}: with $g_\pi = 0$ and $g_h$ re-tuned so memory alone reproduces the
  recorded previous-side AUC, $\Rc(I_p) = 0.078\,[0.019, 0.132]$, positive in 24 of 32 seeds (12 of 16 tasks,
  $p = 0.038$): 49 to 72\% of the recorded score depending on the summary, as Section~\ref{sec:leak} predicts. Reading B (memory
  at calibrated strength) is negative and the null is calibrated (13 of 32).
\item \emph{The history-controlled score}: $R^{2,\perp}_{\mathrm{c}}(I_p) = 0.026$ with a genuine prior and no drift (against 0.170
  uncontrolled, consistent with $\Var(\pi^\perp) \approx 0.19\Var(\pi)$), $-0.017$ for memory alone. In the
  recordings $R^{2,\perp}_{\mathrm{c}}(I_p) = -0.061$ (session median), a median that occurs with probability 0.002 to 0.01 among
  medians of 16 simulated sessions in every condition at calibrated drift strength (except per-unit, 0.015 to 0.12).
  The recordings do something the simulator does not.
\end{enumerate}

\section{The latent state-space model (A18 part B)}
\label{sec:lds}

\paragraph{Model.} One observation per trial, $\bm y_k$ = z-scored $\sqrt{\text{counts}}$ of the pre-stimulus window,
\begin{align}
\bm y_k &= C \bm x_k + D \bm u_k + \bm d + \bm\varepsilon_k, & \bm\varepsilon_k &\sim \mathcal N(0, \operatorname{diag}\bm r), \\
x_{k+1,i} &= \phi_i\, x_{k,i} + \eta_{k,i}, & \eta_{k,i} &\sim \mathcal N(0, 1 - \phi_i^2),\quad \phi_i \in [0, 1 - 10^{-4}],
\end{align}
with $K = 8$ latents, each with its own timescale $\tau_i = -1/\ln\phi_i$ trials, and known inputs $\bm u_k$ = the
cosine drift basis $D_{20}$ at the window time plus $c_{k-1}$, previous choice and previous reward. This is a
linear-Gaussian analogue of GPFA with an exponential (Ornstein--Uhlenbeck) kernel and an explicit drift term. Fitted
by EM (Kalman filter in information form, Rauch--Tung--Striebel smoother; closed-form M-step for $C, D, \bm d, \bm r$;
a bounded 1-D update per $\phi_i$), 3 restarts.

\paragraph{Readout.} To stop the dynamics from \emph{constructing} a prior by integrating past stimuli, the prior is
decoded from each trial's own generalised-least-squares estimate
\begin{equation}
\hat{\bm x}_k = (C^\top R^{-1} C)^{-1} C^\top R^{-1}(\bm y_k - D\bm u_k - \bm d),
\end{equation}
with the same decoder and pseudo sessions as above, giving $\Rc(\mathrm{LDS})$.

\paragraph{Gate (committed before any result).} At the task level (16 tasks): B1 in low-rank drift $\times1$,
$\Rc(\mathrm{LDS}) > 0$ in $\ge 12$ of 16 tasks and mean $\ge \frac12$ mean $\Rc(I_p)$; B2 in memory alone, $< 12$ of
16 positive and mean $\le \frac14$ of the B1 mean; null calibrated. Stop unless B1 and B2 both pass.

\paragraph{Result.} Task-level means (16 tasks):
\begin{center}
\begin{tabular}{lccc}
\toprule
Condition & $\Rc(\mathrm{LDS})$ & tasks $> 0$ & $\Rc(I_p)$ \\
\midrule
calibrated prior, no drift & 0.026 & 13/16 & 0.170 \\
low-rank drift $\times 1$ & $-0.003$ & 8/16 & 0.167 \\
memory alone & $-0.020$ & 6/16 & 0.078 \\
no prior, no memory & $-0.025$ & 6/16 & $-0.043$ \\
\bottomrule
\end{tabular}
\end{center}
B1 fails, so by the rule fixed in advance the latent-model work stops: it is not run on the recordings and no
variant is tried. B2 passes but is moot: the readout is not fooled by memory, but it does not keep a genuine prior
either. Even without drift it keeps only $0.026/0.170 \approx 15\%$ of the population's prior information. In terms of
the model, $\hat{\bm x}_k$ is the GLS projection after subtracting $D\bm u_k$; since $\bm u_k$ contains $c_{k-1}$ and
$\corr(\pi_k, c_{k-1}) = 0.64$, the subtraction removes part of the prior-related activity (about 42\% of $\Var(\pi)$ is
linear in $c_{k-1}$). That is a plausible contributor, not a tested one. With drift, the fitted timescales are
dominated by slow latents (medians $\tau \approx 43, 31, 21, 9$ trials) that absorb slow structure the cosine inputs do not.

%======================================================================
\part{Synthesis}
%======================================================================

\section{The failure, in one set of equations}
\label{sec:synthesis}

\begin{enumerate}
\item The objective \eqref{eq:rrr} keeps the top-$d$ eigenvectors of $G = \hat Y^\top \hat Y + \lambda B_R^\top B_R$:
  directions of largest \emph{predicted variance}. It never sees $\pi$.
\item The recorded prior has $\rho_\pi(\Delta)^2$ falling from 0.99 to 0.01 by 2 min ($\tau_\pi \approx 12$ trials),
  and a variance share about two orders of magnitude below $S_{\mathrm{low}}$. Under the heuristic \eqref{eq:heur} no
  $\Delta$ favours it over a ramp-like drift.
\item In a simulation calibrated to the recordings, the heuristic's prediction holds where the truth is known: without
  drift $V^{\mathrm{in}}_{60}$ captures 41\% of the population's prior information; with drift of the recorded
  strength it captures none, while $\Rc(I_p)$ is unchanged. This does not reproduce $\Delta_\pi < 0$.
\item The measurement itself is ambiguous: $\Rc$ is positive for any population that encodes recent stimuli,
  because $\pi = f(c_{1:k-1})$ and the null redraws $c$. Controlling for recent stimuli removes about 80\% of $\pi$ too.
  So ``the prior is decodable'' and ``the population represents the belief $\pi$'' are different claims, and the data
  here support only the first.
\item A latent model with an explicit drift term, read out trial by trial, did not rescue the prior in simulation
  (Section~\ref{sec:lds}): under calibrated drift it keeps none, and without drift about 15\%.
\end{enumerate}

\paragraph{What would change these conclusions.} A drift model that reproduces $\Delta_\pi < 0$ and the recordings'
$R^{2,\perp}_{\mathrm{c}} < 0$ would strengthen the mechanism; one that cannot would point to features the simulator lacks (unit-level
nonstationarity, longer stimulus memory). Evidence for a belief representation beyond memory would need information
about $\pi$ that a flexible function of recent stimuli cannot provide, for example the run length $\ell_k$ or the
effect of older stimuli at matched recent history, and that question is a good one to take to the instructor.

\section*{Sources}
Code: \texttt{src/pns/\{transform, prior, rrr, lagged, geometry, decode, stats, analysis, timescales, synth, sim\_a18,
lds\}.py}; scripts 04 to 06, 23 to 27. Results: \texttt{results/stats/d5/}, \texttt{results/a16/}, \texttt{results/a17/},
\texttt{results/a18/}. Every decision and its date: \texttt{AMENDMENTS.md}. Every number with its file:
\texttt{RESULTS.md}.

\end{document}
